\exercisegroup

\begin{enumerate}
\def\labelenumi{\arabic{enumi})}
\tightlist
\item
  \begin{enumerate}
  \def\labelenumii{\alph{enumii})}
  \tightlist
  \item
    Let \(\mathbf{E}\) be a locally convex Hausdorff space, and \(\mathbf{E}'\) its dual. Show that the following properties are equivalent:
  \end{enumerate}
\end{enumerate}

α) E is infra-barrelled (III, p.~44, exerc. 7).

β) Every strongly bounded subset in E' is equicontinuous.

\(\gamma)\) Every strongly bounded subset in \(E'\) is relatively weakly compact, and the initial topology of \(E\) is \(\tau(E, E')\) .

δ) The topology induced on E by the strong topology of the bidual \(E''\) is identical with the initial topology on E.

Then a fundamental system of neighbourhoods of 0 for the strong topology of \(\mathbf{E}''\) consists of the closures, for the topology \(\sigma (\mathrm{E}'', \mathrm{E}')\), of a fundamental system of neighbourhoods of 0 for the initial topology of \(\mathbf{E}\) .

\begin{enumerate}
\def\labelenumi{\alph{enumi})}
\setcounter{enumi}{1}
\tightlist
\item
  Prove that if E is infra-barrelled and if its dual \(E'\) is identical with its algebraic dual \(E^{*}\) , then the initial topology of E is the finest locally convex topology.
\end{enumerate}

¶ 2) a) Show that every product of infra-barrelled spaces is infra-barrelled. (Reduce to the case of Hausdorff infra-barrelled spaces; then use exerc. 1 and prop. 15 of IV, p.~14.)

\begin{enumerate}
\def\labelenumi{\alph{enumi})}
\setcounter{enumi}{1}
\tightlist
\item
  Give an example of a Hausdorff and infra-barrelled locally convex space which is neither bornological non barrelled. (Proceed as in III, p.~45, exerc. 16, replacing the barrelled spaces by infra-barrelled spaces and use a).)
\end{enumerate}

\begin{enumerate}
\def\labelenumi{\arabic{enumi})}
\setcounter{enumi}{2}
\item
  Let E be a complex Hausdorff locally convex space, \(E_{0}\) the underlying real locally convex space to E, and let \(E^{\prime}\) and \(E_{0}^{\prime}\) be the duals of E and \(E_{0}\) respectively. Show that the canonical R-linear mapping \(f \mapsto Rf\) from \(E^{\prime}\) onto \(E_{0}^{\prime}\) is a homeomorphism for the S-topologies on \(E^{\prime}\) and \(E_{0}^{\prime}\) , where S is an arbitrary set of bounded subsets of E. From this deduce the definition of the canonical R-linear mapping from the bidual \(E^{\prime\prime}\) onto the bidual \(E_{0}^{\prime\prime}\) , which is a homeomorphism for the weak topologies \(\sigma(E^{\prime\prime}, E^{\prime})\) and \(\sigma(E_{0}^{\prime\prime}, E_{0}^{\prime})\) , as also for the strong topologies \(\beta(E^{\prime\prime}, E^{\prime})\) and \(\beta(E_{0}^{\prime\prime}, E_{0}^{\prime})\) ; by this mapping, E (considered embedded in \(E^{\prime\prime}\) ) transforms into \(E_{0}\) (considered embedded in \(E_{0}^{\prime\prime}\) ).
\item
  Let E be a Hausdorff and infra-barrelled locally convex space.
\end{enumerate}

\begin{enumerate}
\def\labelenumi{\alph{enumi})}
\tightlist
\item
  Show that if the strong dual \(E_{b}^{\prime}\) of E is bornological, then the completion \(\hat{E}\) of E, identified with a vector subspace of \(E^{\prime*}\) (III, p.~21, th. 2), is contained in the bidual \(E^{\prime\prime}\) of E.
\item
  We say that E is distinguished if every subset of \(E^{\prime\prime}\) which is bounded for the topology \(\sigma(\mathrm{E}^{\prime\prime}, \mathrm{E})\) is contained in the closure (for this topology) of a bounded subset of E. Show that for E to be distinguished it is necessary and sufficient that its strong dual \(E_{b}^{\prime}\) is barrelled (cf.~IV, p.~15, prop. 3).
\end{enumerate}

\begin{enumerate}
\def\labelenumi{\arabic{enumi})}
\setcounter{enumi}{4}
\tightlist
\item
  Let \(\mathbf{E}\) be a Hausdorff locally convex space, and \(\mathbf{E}'\) its dual.
\end{enumerate}

\begin{enumerate}
\def\labelenumi{\alph{enumi})}
\item
  In order that the strong topology on \(\mathbf{E}'\) be identical with \(\tau (\mathrm{E}',\mathrm{E})\) , it is necessary and sufficient that \(\mathbf{E}\) is semi-reflexive.
\item
  Suppose that E is infra-barrelled. In order that the strong topology on \(E'\) be identical with the topology of compact convergence, it is necessary and sufficient that E be a Montel space.
\end{enumerate}

¶ 6) Let F and G be two vector spaces in separating duality.

\begin{enumerate}
\def\labelenumi{\alph{enumi})}
\tightlist
\item
  Show that the following properties are equivalent :
\end{enumerate}

\(\alpha)\) the space \(\mathbf{F}\) with a topology compatible with the duality between \(\mathbf{F}\) and \(\mathbf{G}\) , is semi-reflexive;

\(\beta)\) the space \(\mathbf{G}\) , with \(\tau (\mathbf{G},\mathbf{F})\) , is barrelled.

\begin{enumerate}
\def\labelenumi{\alph{enumi})}
\setcounter{enumi}{1}
\tightlist
\item
  Show that the following properties are equivalent :
\end{enumerate}

α) the space F, with τ(F, G), is reflexive;

\(\beta)\) the space \(\mathbf{G}\) , with \(\tau (\mathbf{G},\mathbf{F})\) , is reflexive;

γ) F and G are barrelled for τ(F, G) and τ(G, F) respectively.

\begin{enumerate}
\def\labelenumi{\alph{enumi})}
\setcounter{enumi}{2}
\item
  Show that every locally convex Hausdorff space E, with the topology \(\tau(\mathrm{E}, \mathrm{E}')\) is isomorphic to the quotient of a semi-reflexive space F by a closed subspace. (By III, p.~44, exerc. 14, c), \(E'\) endowed with \(\tau(\mathrm{E}', \mathrm{E})\) is isomorphic to a closed vector subspace M of a Hausdorff barrelled space G; take \(F = G'\) endowed with \(\tau(\mathrm{G}', \mathrm{G})\) and use prop. 11 of IV, p.~10.)
\item
  Let A be an infinite set with \(\operatorname{Card}(\mathrm{A}) > \aleph_{1}\) (S, III, §6, exerc. 10). In the product space \(P = R^{A}\) , let \(E_{0}\) denote the everywhere dense subspace consisting of all \(x = (x_{\alpha})_{\alpha \in \mathrm{A}}\) such that \(x_{\alpha} = 0\) except for a countable set of indices; the space \(E_{0}\) is barrelled and so is the subspace E of P generated by \(E_{0}\) and the point \(1_{A}\) in P all whose coordinates are equal to 1 (III, p.~45, exerc. 16). Let B be a bounded subset in \(E_{0}\) , whose closure in P contains \(1_{A}\) , and let J be a subset of A with cardinality \(\aleph_{1}\) ; let \(pr_{J}\) denote the projection from P onto \(R^{J}\) ; show that there exists a subset \(B_{J}\) in B, with cardinality \(\leqslant \aleph_{1}\) , such that the closure of \(\operatorname{pr}_{J}(B_{J})\) in \(R^{J}\) contains \(1_{J}\) ; then the set \(J' \supset J\) of all indices \(\alpha \in A\) such that for at least one point of \(B_{J}\) the coordinate with index \(\alpha\) is \(\neq 0\) , has a cardinality equal to \(\aleph_{1}\) . We define \(J_{0} = J\) and by induction \(J_{n+1} = J'_{n}\) , and put \(H = \bigcup_{n} J_{n}\) , whose cardinality is \(\aleph_{1}\) , and let \(B_{H} = \bigcup_{n} B_{J_{n}}\) ; show that the closure of \(B_{H}\) (and hence also that of B) contains the point \((1_{H}, 0) \in R^{H} \times R^{A-H} = P\) , which does not belong to \(E_{0}\) . Conclude that for every bounded and compact set C in E, \(C \cap E_{0}\) is again closed in E, and that E is not semi-reflexive.
\item
  Deduce from \(d\) that the dual \(\mathrm{E}'\) of \(\mathrm{E}\) (which can be identified with the dual \(\mathrm{P}' = \mathbf{R}^{(\mathbf{A})}\) of \(\mathrm{P}\) ) is not complete for the topology \(\tau (\mathrm{E}',\mathrm{E})\) , in spite of being semi-reflexive (hence quasi-complete) for this topology (consider the linear form on \(\mathrm{E}\) which is equal to 0 on \(\mathrm{E}_0\) and equal to 1 at the point \(1_{\mathrm{A}}\) , and use III, p.~21, th. 2).
\end{enumerate}

\begin{enumerate}
\def\labelenumi{\arabic{enumi})}
\setcounter{enumi}{6}
\item
  Let \((\mathrm{E}_{i})_{i\in\mathbb{I}}\) be a family of Hausdorff locally convex spaces, P the product space of the \(E_{i}\) , and S their topological direct sum. Show that for P or S to be semi-reflexive (resp. reflexive), it is necessary and sufficient that each \(E_{i}\) be semi-reflexive (resp. reflexive).
\item
  Let E be a Hausdorff locally convex space which is the strict inductive limit of an increasing sequence \((E_{n})\) of closed vector subspaces (II, p.~32, prop. 9).
\end{enumerate}

\begin{enumerate}
\def\labelenumi{\alph{enumi})}
\item
  Show that, if the strong dual of each of the \(\mathbf{E}_n\) is complete, then the strong dual of \(\mathbf{E}\) is complete (III, p.~20, th. 1).
\item
  In order that E be semi-reflexive (resp. reflexive), it is necessary and sufficient that each of the \(E_{n}\) be semi-reflexive (resp. reflexive).
\end{enumerate}

\begin{enumerate}
\def\labelenumi{\arabic{enumi})}
\setcounter{enumi}{8}
\item
  Let \((\mathrm{E}_{\alpha})_{\alpha\in\mathrm{A}}\) be a family of Hausdorff locally convex spaces contained in the same vector space, which is directed for the relation \(\supset\) , such that, if \(E_{\beta} \subset E_{\alpha}\) , the topology of \(E_{\beta}\) is finer than the topology on \(E_{\beta}\) induced by that of \(E_{\alpha}\) . Let E be the intersection of the \(E_{\alpha}\) , endowed with a topology which is the supremum of the topologies on E induced by those on the \(E_{\alpha}\) . Show that if each \(E_{\alpha}\) is semi-reflexive, then E is semi-reflexive (consider an ultrafilter on a bounded subset of E).
\item
  Show that every product, and every topological direct sum of Montel spaces is a Montel space \(^{1}\) .
\item
  Let E be a Hausdorff locally convex space such that the strong dual \(E_{b}^{\prime}\) of E is semi-reflexive. a) Show that, on every strongly bounded subset of \(E^{\prime}\) , the topologies induced by \(\sigma(E^{\prime}, E^{\prime\prime})\) and \(\sigma(E^{\prime}, E)\) are identical.
\end{enumerate}

\begin{enumerate}
\def\labelenumi{\alph{enumi})}
\setcounter{enumi}{1}
\tightlist
\item
  Deduce from a) that E is infra-barrelled for the topology \(\tau(E, E')\) (cf.~IV, p.~52, exerc. 1) and that, if \(\hat{E}\) is its completion for this topology, and is identified with a subset of \(E'^{*}\) (III, p.~21, th. 2), then \(E'' \subset \hat{E}\) . In particular, if E is quasi-complete for \(\tau(E, E')\) , then E is reflexive for this topology.
\end{enumerate}

\(^{1}\) On the other hand, a closed subspace of a Montel space need not be infrabarrelled, and the quotient of a Montel space by a closed subspace need not be semi-reflexive (IV, p.~63, exerc. 8).

\begin{enumerate}
\def\labelenumi{\arabic{enumi})}
\setcounter{enumi}{11}
\tightlist
\item
  Let \(\mathbf{E}\) be a Banach space and \(\mathbf{E}'\) its dual.
\end{enumerate}

\begin{enumerate}
\def\labelenumi{\alph{enumi})}
\item
  Show that the distance \(x \mapsto d(x, \mathbf{A})\) from a point \(x \in \mathbf{E}\) to a closed convex set \(\mathbf{A}\) is a lower semi-continuous function on \(\mathbf{E}\) for the topology \(\sigma(\mathbf{E}, \mathbf{E}')\) .
\item
  Show that if \(\mathrm{E}\) is reflexive, then for every closed convex subset \(\mathrm{A}\) of \(\mathrm{E}\) , there exists a point \(x_0 \in \mathrm{A}\) such that \(\| x_0 \|\) is equal to the distance of 0 to \(\mathrm{A}\) (use \(a\) ). This point is unique if every boundary point of the unit ball of \(\mathrm{E}\) is extremal (II, p.~54, def. 1).
\item
  Suppose E is reflexive and let B be a closed convex and bounded subset in E; deduce from a) and b) that there exist two points \(x \in \mathbf{A}\) , \(y \in \mathbf{B}\) such that \(\| x - y \| = d(\mathbf{A}, \mathbf{B})\) (cf.~V, p.~71, exerc. 8).
\end{enumerate}

\begin{enumerate}
\def\labelenumi{\arabic{enumi})}
\setcounter{enumi}{12}
\item
  Let \(\mathbf{E}\) be a Banach space, and \(\mathbf{M}\) a closed vector subspace of \(\mathbf{E}\) . Show that if \(\mathbf{M}\) and \(\mathrm{E / M}\) are reflexive, then \(\mathbf{E}\) is reflexive.
\item
  Let \(\mathbf{A}\) be an infinite set.
\end{enumerate}

\begin{enumerate}
\def\labelenumi{\alph{enumi})}
\tightlist
\item
  Show that the strong dual of the Banach space \(E = \overline{\mathcal{K}(A)}\) (IV, p.~47, exerc. 1) can be identified with the Banach space \(\ell^{1}(A)\) , and that the strong dual of the Banach space \(\ell'(A)\) can be identified with the Banach space \(\mathcal{B}(A) = \ell^{\infty}(A)\) ; deduce that E is not reflexive and that \(E''/E\) is infinite dimensional (cf.~IV, p.~71, exerc. 18). If A = N, then E and \(E'\) are Banach spaces which satisfy the first axiom of countability, but \(E''\) does not (I, p.~25, exerc. 1).
\item
  Let \(B''\) be the unit ball in \(E'' = \ell^{\infty}(A)\) , and let \(B_{0}''\) be the convex set \(B'' + (B'' \cap E)\) . Show that \(B_{0}''\) is a bounded closed convex set in \(E''\) for the strong topology, with a non-empty interior, but does not have any extremal point. If p is the gauge of \(B_{0}''\) , show that \(E''\) endowed with the norm p is not isometric to any dual of a Banach space and that \(B_{0}''\) is not closed for the topology \(\sigma(E'', E')\) (although \(B''\) is compact for \(\sigma(E'', E')\) and \(B'' \cap E\) is strongly closed).
\end{enumerate}

¶ 15) The notations are those of exerc. 14.

\begin{enumerate}
\def\labelenumi{\alph{enumi})}
\tightlist
\item
  Let \((x_n')\) be a sequence in \(E'\) which converges to 0 for the topology \(\sigma(E', E'')\) ; show that, for every \(\varepsilon > 0\) , there exists a finite subset \(H\) of \(A\) such that \(\sum_{\alpha \notin H} |x_n'(\alpha)| \leqslant \varepsilon\) for every integer \(n\) .
\end{enumerate}

(Argue by reducto ad absurdum: if the property were not true, show that then there exists a number \(\delta > 0\) , an increasing sequence \((n_k)\) of integers, an increasing sequence \((\mathrm{H}_k)\) of finite subsets of A, such that \(\sum_{\alpha \in \mathrm{H}_{k-1}} |x_n'(\alpha)| \leqslant \frac{\delta}{8}\) for \(n \geqslant n_k\) , \(\sum_{\alpha \notin \mathrm{H}_k} |x_n'(\alpha)| \leqslant \frac{\delta}{8}\) for \(n \leqslant n_k\) and

\(\sum_{\alpha \in \mathrm{H}_k - \mathrm{H}_{k-1}} |x_n'(\alpha)| \geqslant \frac{\delta}{2}\) ; show that this implies a contradiction (the « gliding bump » method).)

Deduce that the sequence \((x_{n}^{\prime})\) converges to 0 for the strong topology, although the latter is strictly finer than the topology \(\sigma (\mathrm{E}',\mathrm{E}'')\) .

\begin{enumerate}
\def\labelenumi{\alph{enumi})}
\setcounter{enumi}{1}
\tightlist
\item
  Show that if \((x_{n}^{\prime})\) is a Cauchy sequence in \(E'\) for the topology \(\sigma(E', E'')\) , it converges to a point in \(E'\) for this topology; in other words, \(E'\) is semi-complete (III, p.~7) for \(\sigma(E', E'')\) . (Show that, for every \(\varepsilon > 0\) , there exists a finite subset \(H\) in \(A\) such that \(\sum_{\alpha \notin H} |x_{n}'(\alpha)| \leqslant \varepsilon\) for every integer n; argue by reducto ad absurdum, as in a), and use a).
\end{enumerate}

¶ 16) With the notations of exerc. 14, let E'' be the dual of \(E'' = \ell^{\infty}(A)\).

\begin{enumerate}
\def\labelenumi{\alph{enumi})}
\item
  Let \(e_{\alpha}\) (for \(\alpha \in A\) ) be the element of \(E''\) for which \(e_{\alpha}(\beta) = \delta_{\alpha \beta}\) (Kronecker's symbol). Let \((\mathbf{K}_n)\) be a sequence of finite subsets of \(A\) , two by two disjoint, and let \((x_n''')\) be a sequence of points in \(E'''\) . Show that there exists \((n_k)\) , a strictly increasing infinite sequence of integers \(>0\) such that, if we put \(B = \bigcup_k K_{n_k}\) , then the elements \(y_n' = (x_n'''(\alpha))_{\alpha \in B}\) belong to \(\ell^1(B)\) . (Let \(\delta\) be an arbitrary number \(>0\) , and let \((\mathrm{J}_m)_{m \in \mathbb{N}}\) be a partition of \(\mathbf{N}\) in finite sets. Arguing by reducto ad absurdum, show that if \(x''' \in E'''\) , then there exists an integer \(m\) such that \(|\langle x'', x''' \rangle| \leqslant \delta\) for all \(x'' \in E''\) for which \(\| x'' \| \leqslant 1\) and \(x''(\alpha) = 0\) except for indices \(\alpha\) belonging to the set \(\bigcup_{n \in J_m} K_n\) . Apply this result successively to \(x_1'', x_2'', ...\) in a suitable way.)
\item
  Deduce from \(a\) and from exerc. 15, a) that, if \((x_{n}^{\prime \prime \prime})\) is a sequence converging to 0 in \(\mathrm{E}'''\) for the topology \(\sigma (\mathrm{E}''',\mathrm{E}'')\) and if \(\tilde{x}_n^{\prime \prime \prime}\) is the restriction of \(x_{n}^{\prime \prime \prime}\) to the strongly closed subspace E of \(\mathrm{E}''\) , then \(\lim_{n\to \infty}\| \tilde{x}_n^{\prime \prime \prime}\| = 0\) in \(\mathrm{E}'\) .
\item
  Deduce from b) that, the strongly closed subspace E of \(E''\) does not have a topological complement for the strong topology. (Restricting to the case A = N, let \((e_{n}')\) be the sequence of continuous linear forms on E such that \(\langle x, e_{n}' \rangle = x(n)\) for all \(x \in E\) ; show that the sequence \((e_{n}')\) tends to 0 for \(\sigma(\mathrm{E}', \mathrm{E})\) , but that \(e_{n}'\) cannot be extended to a continuous linear form \(x_{n}'''\) on \(E''\) in such a way that the sequence \((x_{n}'''')\) tends to 0 for \(\sigma(\mathrm{E}'''', \mathrm{E}'')\) .)
\end{enumerate}

\begin{enumerate}
\def\labelenumi{\arabic{enumi})}
\setcounter{enumi}{16}
\tightlist
\item
  Let E be a non-reflexive Banach space, \(E'\) its strong dual, \(E''\) the strong dual of \(E'\) , \(E'''\) the strong dual of \(E''\) and \(E^{IV}\) the strong dual of \(E'''\) .
\end{enumerate}

\begin{enumerate}
\def\labelenumi{\alph{enumi})}
\item
  Show that, in \(E^{\prime\prime}\) , \(E^{\prime}\) and the subspace \(E^{\circ}\) orthogonal to E (when E is considered as a subspace of \(E^{\prime\prime}\) ) are topological complements, and that the projection from \(E^{\prime\prime}\) onto \(E^{\prime}\) for this decomposition is a continuous linear mapping of norm 1. Comparing with exerc. 16, c), deduce that the Banach space \(\mathcal{K}(A)\) is not isomorphic (as a topological vector space) to the strong dual of any Banach space.
\item
  Show that \(E^{IV}\) is the topological direct sum of \(E^{\circ}\) and \(E''\) , and also of \(E^{\prime\circ}\) and \(E^{\circ\circ}\) ; we have \(E'' \cap E^{\circ\circ} = E\) . Let v be the linear mapping from \(E^{IV}\) onto itself which is identity on \(E^{\circ}\) , and on \(E''\) , is the projection from \(E''\) onto \(E^{\circ\circ}\) parallel to \(E^{\circ}\) ; show that v is an isometry, but is not continuous for the topology \(\sigma(\mathrm{E}^{\mathrm{IV}}, \mathrm{E}''')\) .
\end{enumerate}

\begin{enumerate}
\def\labelenumi{\arabic{enumi})}
\setcounter{enumi}{17}
\item
  Show that a Banach space E whose strong dual \(E'\) satisfies the first axiom of countability, and which is semi-complete (III, p.~7) for the weakened topology \(\sigma(\mathrm{E}, \mathrm{E}')\) , is reflexive (compare with IV, p.~54, exerc. 15, b)).
\item
  Let \(\mathrm{E}\) be a Banach space, \(\mathrm{E}'\) its strong dual, \(\mathrm{G}'\) a strongly closed subspace of \(\mathrm{E}'\) satisfying the first axiom of countability for the strong topology. Show that there exists a countable subset of \(\mathrm{E}\) such that, if \(\mathrm{F}\) is the closed vector subspace of \(\mathrm{E}\) generated by this subset, then \(\mathrm{G}'\) is isometric to a strongly closed subspace of the strong dual \(\mathrm{F}'\) of \(\mathrm{F}\) . (Suppose that the sequence \((x_n')\) is strongly dense in \(\mathrm{G}'\) ; for every \(n\) , let \(x_n \in \mathrm{E}\) such that \(\| x_n \| \leqslant 1\) and \(\langle x_n, x_n' \rangle = \left(1 - \frac{1}{n}\right) \| x_n' \|\) ; show that the strongly closed subspace \(F\) of \(E\) generated by the
\end{enumerate}

\[
x _ {n}
\]

is the required space.)

¶ 20) Let E be a Banach space, E' its dual, and B the unit ball in E. In order that every point of B have a countable fundamental system of neighbourhoods for the topology induced by the weakened topology σ(E, E') on B, it is necessary and sufficient that E' satisfies the first axiom of countability for the strong topology. (To show that the condition is necessary, observe that if every point in B has a countable fundamental system of neighbourhoods for the weakened topology, then this is also true for the closure B°° of B in E'' for the topology σ(E'', E'). Therefore there exists a sequence (a'\_n) in E' such that every neighbourhood of 0 in B°° for σ(E'', E') contains the intersection of B°° and of a finite number of polars \{a'\_n\}°; consider the strongly closed subspace W' of E' generated by the a'\_n, and the orthogonal \(W'^{\circ}\) of W' in E''.)

\begin{enumerate}
\def\labelenumi{\arabic{enumi})}
\setcounter{enumi}{20}
\tightlist
\item
  Let E be a Banach space, \(E'\) its dual, B the unit ball in E and \(B_{r}'\) the closed ball with centre 0 and radius r in \(E'\) .
\end{enumerate}

\begin{enumerate}
\def\labelenumi{\alph{enumi})}
\item
  Let \(M_{1}^{\prime}\) , \(M_{2}^{\prime}\) be two vector subspaces of \(E^{\prime}\) , which are everywhere dense for the weak topology \(\sigma(\mathrm{E}^{\prime}, \mathrm{E})\) . In order that the topologies induced by \(\sigma(\mathrm{E}, \mathrm{M}_{1}^{\prime})\) and \(\sigma(\mathrm{E}, \mathrm{M}_{2}^{\prime})\) on B coincide, it is necessary and sufficient that the strong closures of \(M_{1}^{\prime}\) and \(M_{2}^{\prime}\) in \(E^{\prime}\) are identical.
\item
  Let \(M'\) be a vector subspace of \(E'\) which is everywhere dense for \(\sigma(E', E)\) ; let \(\mathbf{M}'^{(1)}\) denote the vector subspace generated by the closure of \(M' \cap B_{1}'\) in \(E'\) for the topology \(\sigma(E', E)\) . In order that \(\mathbf{M}'^{(1)} = E'\) , it is necessary and sufficient that the weak closure of \(M' \cap B_{1}'\) contains a ball \(B_{r}'\) with r \textgreater{} 0 (use the fact that \(E'\) is barrelled for the strong topology).
\item
  Let \(r\) be the supremum of the numbers \(t\) such that the weak closure of \(\mathbf{M}' \cap \mathbf{B}_1'\) contains a ball \(\mathbf{B}_t'\) ; the number \(r\) is said to be the characteristic of \(\mathbf{M}'\) . Show that \(r\) is the infimum of the numbers \(\sup_{x' \in \mathbf{M}' \cap \mathbf{B}_1'} |\langle x, x' \rangle| / \| x\|\) where \(x\) ranges over the set of all points \(\neq 0\) in \(E\) (use the Hahn-Banach theorem).
\item
  Show that \(1 / r\) is the supremum of \(\| x\|\) as \(x\) ranges over the closure of B in E for the topology \(\sigma (\mathrm{E},\mathbf{M}^{\prime})\) (use \(c\) ) and the Hahn-Banach theorem).
\item
  Let \(M^{\prime\circ}\) be the orthogonal of \(M^{\prime}\) in \(E^{\prime\prime}\) ; show that \(r = \inf(\|x + z^{\prime\prime}\|/\|x\|)\) where \(z^{\prime\prime}\) ranges over \(M^{\prime\circ}\) and x ranges over the set of non-zero points in E (use c) and the Hahn-Banach th). Deduce that in order that \(\mathbf{M}^{\prime(1)} = \mathbf{E}^{\prime}\) , it is necessary and sufficient that \(E + M^{\prime\circ}\) be strongly closed in \(E^{\prime\prime}\) (use Banach's th., I, p.~17).
\end{enumerate}

\(f\) ) Let \(\mathbf{A} = \mathbf{N}\times \mathbf{N}\) and \(\mathrm{E} = \overline{\mathcal{K}(\mathbf{A})}\) (cf.~IV, p.~54, exerc. 14). In the space \(\mathrm{E}'' = \ell^{\infty}(\mathrm{A})\) , let \(\mathrm{P}\) be the vector subspace consisting of all points \(x = (x_{ij})\) such that \(x_{ij} = x_{0j} / (j + 1)\) for all \(i\geqslant 0\) . Show that \(\mathbf{P} = \mathbf{M}'^{\circ}\) where \(\mathbf{M}'\) is a vector subspace of \(\mathrm{E}'\) that is everywhere dense (for \(\sigma (\mathrm{E}',\mathrm{E}))\) , but that \(\mathrm{E} + \mathrm{M}'^{\circ} = \mathrm{E} + \mathrm{P}\) is not strongly closed in \(\mathrm{E}''\) ; deduce that the characteristic of \(\mathbf{M}'\) is 0.

\begin{enumerate}
\def\labelenumi{\arabic{enumi})}
\setcounter{enumi}{21}
\tightlist
\item
  Let E be a Banach space, \(E'\) its dual and \(M'\) a strongly closed subspace in \(E'\) which is everywhere dense for the weak topology on \(E'\) . We say that M is irreducible if there exists no vector subspace \(N' \neq M'\) of \(M'\) which is strongly closed and weakly everywhere dense in \(E'\) . a) Show that \(M'\) is irreducible if and only if the orthogonal \(M'^{\circ}\) of \(M'\) in \(E''\) is the topological complement of E (for the strong topology of \(E''\) ). Deduce that then \(\mathbf{M}^{(1)} = \mathbf{E}'\) (exerc. 21) and that E is isomorphic to the strong dual of the space \(M'\) endowed with the topology induced by the strong topology of \(E'\) .
\end{enumerate}

\begin{enumerate}
\def\labelenumi{\alph{enumi})}
\setcounter{enumi}{1}
\item
  Show that \(\mathbf{M}'\) is irreducible if and only if the unit ball in \(\mathbf{E}\) is relatively compact for the topology \(\sigma (\mathrm{E},\mathbf{M}')\) (use exerc. 21, a)).
\item
  For E to be isomorphic to a strong dual of a Banach space (for the topological vector space structure), it is necessary and sufficient that there exist an irreducible subspace in \(E'\) . Deduce a new proof of the fact that the Banach space \(\overline{\mathcal{K}(\mathbf{N})}\) is not isomorphic to a strong dual of a Banach space (cf.~IV, p.~55, exerc. 16, c)).
\end{enumerate}

\begin{enumerate}
\def\labelenumi{\arabic{enumi})}
\setcounter{enumi}{22}
\tightlist
\item
  With the same notations as in exerc. 22, assume that \(\mathbf{M}'\) is irreducible.
\end{enumerate}

\begin{enumerate}
\def\labelenumi{\alph{enumi})}
\item
  In order that the canonical mapping from E into \(E^{\prime\prime}/M'^{\circ}\) , which is the restriction of the canonical mapping \(E^{\prime\prime} \rightarrow E^{\prime\prime}/M'^{\circ}\) be a Banach space isometry, it is necessary and sufficient that the characteristic (IV, p.~55, exerc. 21) of \(M^{\prime}\) is equal to 1. Then, \(M^{\prime}\) endowed with the norm induced by that of \(E^{\prime}\) is said to be the predual of E and E can be canonically identified (with its norm) with the dual of the Banach space \(M^{\prime}\) .
\item
  For every vector subspace F of E which is closed for \(\sigma (\mathrm{E},\mathbf{M}^{\prime})\) , show that the canonical image of \(\mathbf{M}'\) in the dual \(\mathrm{F}'\) of F, identified with \(\mathrm{E}^{\prime} / \mathrm{F}^{\circ}\) , is a predual of F.
\end{enumerate}

\begin{enumerate}
\def\labelenumi{\arabic{enumi})}
\setcounter{enumi}{23}
\tightlist
\item
  \begin{enumerate}
  \def\labelenumii{\alph{enumii})}
  \tightlist
  \item
    Let \((a_{k})_{1\leqslant k\leqslant n}\) be a finite sequence of points in a normed space \(E\) , and let \((\lambda_k)_{1\leqslant k\leqslant n}\) be a finite sequence of numbers \(>0\) such that \(\sum_{k = 1}^{n}\lambda_{k} < 1\) ; put \(\mu_{k} = 1 - \sum_{j = 1}^{k - 1}\lambda_{j}\) for every \(k\) ; then
  \end{enumerate}
\end{enumerate}

\[
\left\| \sum_ {k = 1} ^ {n} \lambda_ {k} a _ {k} \right\| \leqslant \frac {\lambda_ {n}}{\mu_ {n - 1}} \left\| \mu_ {n - 1} a _ {n} + \sum_ {k = 1} ^ {n - 1} \lambda_ {k} a _ {k} \right\| + \frac {\mu_ {n}}{\mu_ {n - 1}} \left\| \sum_ {k = 1} ^ {n - 1} \lambda_ {k} a _ {k} \right\|.
\]

\begin{enumerate}
\def\labelenumi{\alph{enumi})}
\setcounter{enumi}{1}
\tightlist
\item
  Let \((a_{n})\) be an infinite sequence of points in the unit ball of E, and let \((\lambda_{n})\) be an infinite sequence of numbers \(>0\) such that \(\sum_{n} \lambda_{n} = 1\) . For every \(n > 0\) , put \(\mu_{n} = 1 - \sum_{k=1}^{n-1} \lambda_{k}\) and \(b_{n} = \sum_{k=1}^{n-1} \lambda_{k} a_{k} + \mu_{n} a_{n}\) ; show that, for all \(n \geqslant 1\) , we have
\end{enumerate}

\[
\left\| \sum_ {k = 1} ^ {n} \lambda_ {k} a _ {k} \right\| \leqslant \mu_ {n} \sum_ {k = 1} ^ {n} \frac {\lambda_ {k} \| b _ {k} \|}{\mu_ {k - 1} \mu_ {k}}.
\]

(Apply a) by induction.)

\begin{enumerate}
\def\labelenumi{\alph{enumi})}
\setcounter{enumi}{2}
\tightlist
\item
  Let \((\mathrm{C}_n)\) be a decreasing sequence of convex sets in E, contained in the unit ball, and suppose that \(d(0,\mathrm{C}_1)\geqslant \theta >0\) (hence, a fortiori \(d(0,\mathrm{C}_n)\geqslant \theta\) for all \(n\) ). Let \((\lambda_{n})\) be a sequence of numbers \(>0\) such that \(\sum_{n}\lambda_{n} = 1\) . Show that there exists a number \(\alpha\) such that \(\theta \leqslant \alpha \leqslant 1\) and a sequence \((x_{n})\) of points of E such that \(x_{n}\in \mathbf{C}_{n}\) for all \(n\) , \(\| \sum_{n}\lambda_{n}x_{n}\| = \alpha\) and, for all \(n\)
\end{enumerate}

\[
\left\| \sum_ {k = 1} ^ {n} \lambda_ {k} x _ {k} \right\| \leqslant \alpha (1 - \theta \sum_ {j = n + 1} ^ {\infty} \lambda_ {j}).
\]

(Take \(x_{1}\) such that \(\| x_{1}\|\) is arbitrarily close to \(d(0, \mathrm{C}_{1})\) , then, by induction, take \(x_{n}\) such that \(\left\| \sum_{k=1}^{n-1} \lambda_{k} x_{k} + \left(\sum_{j=n}^{\infty} \lambda_{j}\right) x_{n}\right\|\) is arbitrarily close to the infimum of the numbers \(\left\| \sum_{k=1}^{n-1} \lambda_{k} x_{k} + \left(\sum_{j=n}^{\infty} \lambda_{j}\right) y\right\|\) where \(y\) ranges over \(\mathrm{C}_n\) . Then use \(b)\) .)

¶ 25) Let E be a Banach space satisfying the first axiom of countability. Show that the following properties are equivalent :

α) E is not reflexive.

\(\beta)\) For every number \(\theta\) such that \(0 < \theta < 1\) , there exists a sequence \((x_n')\) in \(\mathbf{E}'\) such that, \(\| x_n'\| \leqslant 1\) for all \(n\) , that the sequence \((x_n')\) converges to 0 for \(\sigma (\mathrm{E}',\mathrm{E})\) and that the distance of 0 from the convex set generated by the \(x_{n}^{\prime}\) is \(\geqslant \theta\) .

\(\gamma)\) For every number \(\theta\) such that \(0 < \theta < 1\) and every sequence \(\lambda_{n}\) of numbers \(>0\) such that \(\sum_{n} \lambda_{n} = 1\) , there exists a number \(\alpha\) such that \(\theta \leqslant \alpha \leqslant 1\) and a sequence \((y_{n}')\) of points of

E' such that \(\| y_n'\| \leqslant 1\) for all \(n\) , \(\left\| \sum_{n}\lambda_ny_n'\right\| = \alpha\) and \(\left\| \sum_{k=1}^{n}\lambda_ky_k'\right\| \leqslant \alpha (1 - \theta \sum_{j=n+1}^{\infty}\lambda_j)\) for all \(n\) .

\(\delta)\) There exists \(z' \in \mathrm{E}'\) such that for no \(x \in \mathrm{E}\) do we have \(\left|\langle x, z' \rangle\right| = \| x \| \cdot \| z'\|\) (Theorem of James-Klee).

(To see that \(\alpha\) ) implies \(\beta\) ), observe that there exists \(z'' \in \mathrm{E}''\) such that \(\| z'' \| < 1\) and \(d(z'', \mathrm{E}) > \theta\) . If \((x_n)\) is an everywhere dense sequence in \(\mathrm{E}\) , find the sequence \((x_n')\) in \(\mathrm{E}'\) such that \(\| x_n' \| < 1\) and such that \(\langle x_k, x_n' \rangle = 0\) for \(k \leqslant n\) and \(\langle x_n', z'' \rangle = \theta\) . To see that \(\beta\) ) implies \(\gamma\) ), use exerc. 24. For \(\gamma\) ) implies \(\delta\) ), show that for all \(x \in \mathrm{E}\) we have \(|\sum_{n} \lambda_n \langle x, y_n' \rangle| < \alpha\) , with the notations of \(\gamma\) ).

\begin{enumerate}
\def\labelenumi{\arabic{enumi})}
\setcounter{enumi}{25}
\tightlist
\item
  A locally convex space E is said to have the property (GDF) if every linear mapping u from E into a Banach space F which satisfies the following property, is continuous: in the product space \(E \times F\) , every limit of a convergent sequence of points in the graph \(\Gamma\) of u again belongs to \(\Gamma\) . Every Fréchet space has the property (GDF) (I, p.~19, cor. 5); this is also true of every inductive limit of a family of Fréchet spaces (II, p.~34, prop. 10). Show that every Hausdorff locally convex space with the (GDF) property is barrelled. (Let V be a barrel in E, q its gauge and H the Hausdorff space associated with E endowed with this semi-norm; show that the canonical mapping \(\pi\) from E into the completion \(\hat{H}\) is continuous by using the property (GDF) and the fact that every linear form \(x' \in V^{0}\) can be extended uniquely to a continuous linear form on \(\hat{H}\) , the set of these forms being the unit ball in the dual of \(\hat{H}\) .)
\end{enumerate}

